% Options for packages loaded elsewhere
\PassOptionsToPackage{unicode}{hyperref}
\PassOptionsToPackage{hyphens}{url}
\documentclass[]{article}
\usepackage{xcolor}
\usepackage{amsmath,amssymb}
\setcounter{secnumdepth}{-\maxdimen}
\usepackage{iftex}
\ifPDFTeX
  \usepackage[T1]{fontenc}
  \usepackage[utf8]{inputenc}
  \usepackage{textcomp}
\else
  \usepackage{unicode-math}
  \defaultfontfeatures{Scale=MatchLowercase}
  \defaultfontfeatures[\rmfamily]{Ligatures=TeX,Scale=1}
\fi
\usepackage{lmodern}
\IfFileExists{upquote.sty}{\usepackage{upquote}}{}
\IfFileExists{microtype.sty}{\usepackage[]{microtype}\UseMicrotypeSet[protrusion]{basicmath}}{}
\makeatletter
\@ifundefined{KOMAClassName}{
  \IfFileExists{parskip.sty}{
    \usepackage{parskip}
  }{
    \setlength{\parindent}{0pt}
    \setlength{\parskip}{6pt plus 2pt minus 1pt}
  }
}{
  \KOMAoptions{parskip=half}
}
\makeatother
\usepackage{longtable,booktabs,array}
\newcounter{none}
\usepackage{calc}
\usepackage{etoolbox}
\makeatletter
\patchcmd\longtable{\par}{\if@noskipsec\mbox{}\fi\par}{}
\makeatother
\IfFileExists{footnotehyper.sty}{\usepackage{footnotehyper}}{\usepackage{footnote}}
\makesavenoteenv{longtable}
\setlength{\emergencystretch}{3em}
\providecommand{\tightlist}{\setlength{\itemsep}{0pt}\setlength{\parskip}{0pt}}
\usepackage{bookmark}
\IfFileExists{xurl.sty}{\usepackage{xurl}}{}
\urlstyle{same}
\hypersetup{hidelinks,pdfcreator={LaTeX via pandoc}}
\usepackage[a4paper,margin=2.5cm]{geometry}
\author{}
\date{}

\begin{document}

\section{Assignment 2 -- Formative assessment (Part 2)}
\label{assignment-2-formative-assessment-part-2}

\textbf{Question 3 and the plan of numerical experiments of Question 2}

\textbf{Group:} 16
\textbf{Members:} William Ammentorp -- s223860, Anton H{\o}ier Troelsen -- s223877,
Emil M{\o}ller Kristensen -- s244504, Selim Gabriele Francesco Foglino -- s260275

\emph{At most 6 pages, plus the feedback page. Keep the items, their order and their numbering.}

\begin{center}\rule{0.5\linewidth}{0.5pt}\end{center}

\subsection{Part 0 -- Status and feedback requests}
\label{part-0-status-and-feedback-requests}

\subsubsection{0.1 Status}
\label{status}

For each model: formulated, implemented, validated (yes / partly / not yet), and the validation check used.

\begin{longtable}[]{@{}lllll@{}}
\toprule
Model & Formulated & Implemented & Validated & Validation check used \\
\midrule
\endhead
\(\mathcal{DP}\), \(\mathcal{DP}^{dis}\), \(\mathcal{DP}^{bb}\), and
\(\mathcal{DP}^{disbb}\) & Yes & Yes & Partly & MIP gap 0.00\%, KKT construction,
base-case and sensitivity runs, aggregate identity checks \\
Baseline \(\mathcal{DP}\) at the expected parameters & Yes & Yes & Yes &
Comparison with model \(\mathcal{C}\): equal objective and zero cap exceedance
at \(\alpha=5\) DKK/kWh \\
DP individual-rationality variants & Partly & No & Not yet &
No DP IR constraints or DP member-comparison results exist yet \\
Scenario generation & No & No & Not yet &
Not started \\
\(\mathcal{SDP}^{E}\) & No & No & Not yet &
Not started \\
\(\mathcal{SDP}^{CC}\) & No & No & Not yet &
Not started \\
Evaluation module (out-of-sample) & No & No & Not yet &
Not started \\
\bottomrule
\end{longtable}

The DP implementations use Big-M complementarity in \texttt{add\_kkt}.
They include uniform or member-specific posted prices and an optional
budget-balance constraint. The Q1 \texttt{settle\_Q1} routine evaluates
individual rationality for model \(\mathcal C\), but it is not yet applied
inside or across the DP variants.

If you took part in Round 1, what you changed after the feedback (one line):

\begin{quote}
We did not take part in Round 1. This sheet reports the first formative-assessment plan.
\end{quote}

\subsubsection{0.2 Feedback requests}
\label{feedback-requests}

\begin{enumerate}
\tightlist
\item
  \textbf{D1, IR:} Should individual rationality in DP be imposed against
  model \(\mathcal R\), model \(\mathcal I\), or a best-alternative cost
  \(\min(c_i^R,c_i^I)\) for each member and scenario?
\item
  \textbf{D1, infeasibility:} When a DP budget-balance or IR requirement
  becomes infeasible, should the study report the nearest feasible
  relaxation or the unsolved status as a separate infeasible outcome?
\item
  \textbf{D2:} Is the price bound \(|x_t|\le X\) sufficient, or should
  the feedback suggest a tighter/regularized interpretation of the posted
  price in addition to Big-M bounds?
\end{enumerate}

\begin{center}\rule{0.5\linewidth}{0.5pt}\end{center}

\subsection{Part D -- Question 2, plan of numerical experiments}
\label{part-d-question-2-plan-of-numerical-experiments}

One experiment card per planned study. D1 and D2 are expected; D3 is
optional.

\subsubsection{D1 Three-step study: \(\mathcal{DP}\) with the cap only,
then with budget balance, then with individual rationality, against the
relevant models of Question 1}
\label{d1-three-step-study-mathcaldp}

\begin{longtable}[]{@{}
>{\raggedright\arraybackslash}p{(\linewidth - 2\tabcolsep) * \real{0.50}}
>{\raggedright\arraybackslash}p{(\linewidth - 2\tabcolsep) * \real{0.50}}@{}}
\toprule
Field & \\
\midrule
\endhead
\textbf{1. Question} &
How much of the centralized flexibility and welfare of model
\(\mathcal C\) can the manager recover using only hourly posted prices?
How does adding budget balance and then individual rationality change
efficiency, feasibility, member payments, and the use of discriminatory
prices? \\
\textbf{2. Hypothesis} &
At the base penalty and price bound, the cap-only variants
\(\mathcal{DP}\) and \(\mathcal{DP}^{dis}\) should reproduce much of
\(\mathcal C\). Budget balance may require the community to give up
flexibility if no uniform price can provide both the central dispatch and
zero net budget. Discriminatory prices may reduce this loss, but the
members with high disutility or large PV/battery exports may be the
main blockers of uniform-price IR and budget redistribution. \\
\textbf{3. Protocol} &
Use the base-case community and day \(d1\). Run
\(\mathcal R\), \(\mathcal I\), \(\mathcal C\),
\(\mathcal{DP}\), \(\mathcal{DP}^{dis}\),
\(\mathcal{DP}^{bb}\), and \(\mathcal{DP}^{disbb}\).
The cap-only step uses no \texttt{budget\_balance} constraint. The second
step adds \texttt{budget\_balance == 0}. The planned third step adds
member-level IR constraints versus \(\mathcal R\), \(\mathcal I\), or
their best alternative. Infeasibility must be recorded explicitly with
status, parameters, and the resulting blocking members. \\
\textbf{4. Metrics} &
Cost loss versus \(\mathcal C\), community objective, aggregate and
member shifted energy, import peak, cap exceedance, manager budget,
member-payment range, number of IR violations, and maximum member
price spread. \\
\textbf{5. Figure or table} &
\emph{Figure or table placeholder (required):}
\textbf{A four-panel comparison with cost loss, budget, cap exceedance,
and IR violations across \(\mathcal C\), \(\mathcal R\),
\(\mathcal I\), and the four DP variants.} Use model variants on the
x-axis and one metric per panel. The draft caption should conclude
whether budget balance and IR are achieved without a meaningful loss
from the centralized dispatch. \\
\textbf{6. Expected insight} &
The base result is expected to show that cap-only DP can recover the
centralized physical dispatch; the budget-balance and IR steps reveal
whether redistribution can be obtained without sacrificing efficiency.
A result contradicting this would be an efficiency loss already in the
cap-only problem or, conversely, no loss after adding budget balance. \\
\textbf{7. Status} &
Preliminary: cap-only and both budget-balance variants have been solved.
DP versus model C base-case results have been validated. IR-constrained
DP variants and infeasibility reporting are planned. \\
\bottomrule
\end{longtable}

\begin{figure}[htbp]
\centering
\fbox{\begin{minipage}{0.94\linewidth}
\centering
\vspace{2cm}
\textbf{PICTURE NEEDED: D1 base-case comparison figure}\\
Panels: (1) total cost and loss versus model C, (2) manager budget,
(3) peak import and exceedance, and (4) member-level IR gains/violations
against R and I. Compare R, I, C, DP, DP-dis, DP-bb, and DP-disbb.
\vspace{2cm}
\end{minipage}}
\caption{Placeholder for the D1 three-step study.}
\label{fig:d1_placeholder}
\end{figure}

\begin{longtable}[]{@{}llllll@{}}
\toprule
Metric & Definition & Unit & Reference & Point of view & Sign \\
\midrule
\endhead
Cost loss & \(J(\mathcal M)-J(\mathcal C)\), \(J\) is the community objective
& DKK & \(\mathcal C\) & Manager/community & Positive is worse \\
Manager budget & Member payments minus grid cost and exceedance penalty
& DKK & 0 & Manager & Positive is surplus \\
Member cost gain & Cost in alternative minus cost under the DP variant
& DKK & R or I & Member & Positive is better off \\
IR violations & Number of members with gain \(<0\) & count & 0 & Member
& Higher is worse \\
Peak import & Maximum aggregate connection-point import & kW & Cap profile
& Community/DSO & Lower is better \\
Exceedance & Sum of hourly import above the cap & kWh & 0 & DSO
& Positive is worse \\
Price spread & Maximum difference between member prices in an hour
& DKK/kWh & 0 for uniform DP & Members & Positive indicates discrimination \\
\bottomrule
\end{longtable}

\textbf{Preliminary base-case observation.}
At \(\alpha=5\) DKK/kWh and the default price bound, the community objective
is 360.19 DKK for DP and DP-dis, and also for both budget-balanced variants.
The aggregate peak import is 9.39 kW and exceedance is zero. The
budget-balanced runs have approximately zero manager budget, while the
cap-only runs retain a positive budget. These numbers show that a
uniform/common or discriminatory price can affect payments while reproducing
the same base-case physical dispatch.

\subsubsection{D2 Contract terms: whether a uniform price can meet the
requirements, with the feasibility of each step, its efficiency cost,
and the blocking members along the study}
\label{d2-contract-terms}

\begin{longtable}[]{@{}
>{\raggedright\arraybackslash}p{(\linewidth - 2\tabcolsep) * \real{0.50}}
>{\raggedright\arraybackslash}p{(\linewidth - 2\tabcolsep) * \real{0.50}}@{}}
\toprule
Field & \\
\midrule
\endhead
\textbf{1. Question} &
How do capacity-exceedance penalty \(\alpha\), posted-price bound \(X\),
and cap tightness affect uniform-price feasibility, efficiency, and
budget balance relative to the centralized and discriminatory benchmarks? \\
\textbf{2. Hypothesis} &
The uniform models should lose less as \(\alpha\) and the price bound
increase. Very low price bounds and tight caps should be where uniform
prices fail to reproduce model C or become budget-unbalanced.
Discriminatory prices should reduce the loss under such restrictions.
The members that block uniform IR are expected to be members with little
flexibility or with high PV/battery surplus; this must be confirmed once
DP member-level IR constraints are implemented. \\
\textbf{3. Protocol} &
Vary one contract term at a time on the base community:
\(\alpha\in\{0,2.5,5,7.5,10,12.5,15,20\}\) DKK/kWh;
\(X\in\{3,5,7,10,13,16,20\}\) DKK/kWh; and the cap scale
\(\{0.6,0.8,1.0,1.2,1.5\}\), where 1.0 is the shipped base cap.
Keep the remaining terms fixed. Solve model C as a benchmark and
\(\mathcal{DP}\), \(\mathcal{DP}^{dis}\),
\(\mathcal{DP}^{bb}\), and \(\mathcal{DP}^{disbb}\) for each setting.
The sensitivity currently contains 20 parameter points and 100 model
result rows, including model C and the four DP variants. Record infeasible
models separately. \\
\textbf{4. Metrics} &
Cost loss versus C, manager budget, exceedance, hours above cap, peak
import, feasibility status, member-payment range and spread, and the
blocking members under budget/IR requirements once implemented. \\
\textbf{5. Figure or table} &
\emph{Figure or table placeholder (required):}
\textbf{Three sensitivity panels for (a) penalty \(\alpha\),
(b) price bound \(X\), and (c) cap scale.}
Plot cost loss versus C and manager budget, with C, DP, DP-dis,
DP-bb, and DP-disbb identified by separate markers. \\
\textbf{6. Expected insight} &
The uniform-price feasible region and its efficiency cost can be identified
from the sensitivity figures. The discriminatory variants should remain
closer to C when uniform prices become restrictive.
A contradicting result would be DP-dis not improving the uniform result
for every restrictive parameter, or a uniform budget-balanced model
remaining exactly at C without a restriction. \\
\textbf{7. Status} &
Preliminary results exist for all three sweeps. The CSV results
provide 20 parameter points and the corresponding C/DP variants.
DP individual-rationality constraints and blocking-member analysis
are not yet implemented. \\
\bottomrule
\end{longtable}

\begin{figure}[htbp]
\centering
\fbox{\begin{minipage}{0.94\linewidth}
\centering
\vspace{2cm}
\textbf{PICTURE NEEDED: D2 sensitivity plot}\\
Create \texttt{Q2\_sensitivity.png} already produced by
\texttt{plot\_dp\_sweep}. It should show cost loss versus C and manager
budget for penalty, price bound, and cap scale. Confirm the x-axis and
legend before publication.
\vspace{2cm}
\end{minipage}}
\caption{Placeholder for the D2 contract-term sensitivity.}
\label{fig:d2_placeholder}
\end{figure}

\begin{longtable}[]{@{}llllll@{}}
\toprule
Metric & Definition & Unit & Reference & Point of view & Sign \\
\midrule
\endhead
Cost loss & \(J_{\mathrm{DP}}-J_{\mathrm C}\) & DKK & 0 & Community
& Positive is worse \\
Feasibility & \texttt{OK} or \texttt{NO SOLUTION} & category & OK
& Manager & An infeasible outcome is a restriction \\
Budget & Member payments minus DSO grid cost & DKK & 0 & Manager
& Absolute value measures imbalance \\
Exceedance & Aggregate import above the cap & kWh & 0 & DSO
& Positive is worse \\
Peak import & Maximum aggregate import & kW & Cap & DSO
& Higher than cap indicates congestion \\
Hours above cap & Number of hours with positive exceedance & hours
& 0 & DSO & Positive is worse \\
Member spread & Difference among member-specific posted prices
& DKK/kWh & 0 & Members & Positive indicates discrimination \\
Blocking members & Members that prevent the uniform/IR requirement
& count/IDs & 0 & Members & Higher count is worse \\
\bottomrule
\end{longtable}

\textbf{Preliminary parameter evidence.}
At a penalty of zero, model C and all DP variants have the same high
exceedance of 137.18 kWh and a peak of 43.44 kW. At the base penalty,
DP and DP-dis recover model C exactly, while uniform and discriminatory
budget balance both remain feasible. With a price bound of 3 DKK/kWh,
DP has a 1.239 DKK loss, DP-dis 0.190 DKK, DP-bb 1.555 DKK, and
DP-disbb 0.354 DKK versus C. At a 10 DKK/kWh bound, all DP variants
essentially reproduce C.

\subsubsection{D3 Optional: one further study, and the conclusion it
would change}
\label{d3-optional-one-further-study-and-the-conclusion-it-would-change}

\begin{longtable}[]{@{}ll@{}}
\toprule
Field & \\
\midrule
\midrule
\endhead
\textbf{1. Question} &
Does removing battery availability reduce the ability of uniform posted
prices to reproduce model C and satisfy budget balance? \\
\textbf{2. Hypothesis} &
Without batteries, fewer flexible assets remain for intraday shifting.
Uniform prices are therefore expected to become more restrictive and to
have a larger efficiency loss, while member-specific prices may preserve
more of model C. \\
\textbf{3. Protocol} &
Create a no-battery member-mix variation by setting all battery parameters
to \texttt{None}. Recompute model R and the connection-point cap
according to the Table 11 rule, and recompute individual-contract shares
using the selected split rule. Run C and the four DP variants in the base
day, keeping \(\alpha\), \(\beta\), and \(X\) fixed. \\
\textbf{4. Metrics} &
Cost loss versus C, budget, cap exceedance, shifted energy, number of
uniform-price/IR violations, and the member types blocked. \\
\textbf{5. Figure or table} &
\emph{Figure or table placeholder (required):} paired bars or lines for
base mix versus no-battery mix, showing cost loss and IR violations for
uniform and discriminatory DP variants. \\
\textbf{6. Expected insight} &
If discriminatory DP recovers C in the no-battery community but uniform
DP does not, then batteries are not the only source of discrimination
requirements; load-disutility heterogeneity alone can block uniform prices.
The conclusion would change if both uniform and discriminatory models
remain at C or if removing batteries does not affect the uniform result. \\
\textbf{7. Status} &
Optional and planned. The variation is not yet implemented or solved. \\
\bottomrule
\end{longtable}


